\documentclass[11pt]{article}
\usepackage{amsmath, amssymb, geometry, graphicx, enumerate, xcolor}
\geometry{a4paper, margin=1in}

\title{Homework 16.1 Vector Fields}
\author{Tashfeen Omran}
\date{\today}

\begin{document}

\maketitle

\tableofcontents
\vspace{0.5cm}
\vspace{0.5cm}

% ==============================================================================
% PART 1: Comprehensive Introduction
% ==============================================================================
\section{Part 1: Comprehensive Introduction, Context, and Prerequisites}

\subsection{Core Concepts}
A \textbf{Vector Field} is a function that assigns a vector to every point in a space. Unlike scalar functions $f(x, y)$, which output a single number (a scalar) for each point, a vector field $\mathbf{F}(x,y)$ outputs an arrow (a vector) with both a magnitude and a direction. 

If we are in two-dimensional space $\mathbb{R}^2$, a vector field is written as:
\[ \mathbf{F}(x, y) = P(x, y)\mathbf{i} + Q(x, y)\mathbf{j} = \langle P(x, y), Q(x, y) \rangle \]
Here, $P(x,y)$ and $Q(x,y)$ are ordinary scalar functions that define the horizontal and vertical components of the vector at point $(x,y)$. 

In three-dimensional space $\mathbb{R}^3$, a vector field adds a third component:
\[ \mathbf{F}(x, y, z) = P(x, y, z)\mathbf{i} + Q(x, y, z)\mathbf{j} + R(x, y, z)\mathbf{k} = \langle P, Q, R \rangle \]

Visually, you can imagine a vector field as a room filled with arrows. At every single coordinate in the room, there is an arrow floating. The length and direction of the arrow change depending on where you are standing. 

\subsection{Intuition and Derivation: The Gradient Field}
Where do vector fields come from? One of the most important ways to generate a vector field is by taking the \textbf{gradient} of a scalar function $f(x,y)$. 
Recall that the gradient $\nabla f$ is a vector composed of partial derivatives:
\[ \nabla f(x, y) = \langle f_x(x, y), f_y(x, y) \rangle \]
Because evaluating $\nabla f$ at any point $(x,y)$ gives a vector, $\nabla f$ is itself a vector field! We call this a \textbf{gradient vector field}. 

The intuition is powerful: if $f(x,y)$ represents the elevation of a mountain, the gradient vector field $\nabla f$ tells you which direction to step to go uphill the steepest. The arrows point toward higher elevations, and their length represents the steepness of the slope.

\subsection{Historical Context and Motivation}
Vector fields were developed in the 19th century to mathematically formalize the invisible forces studied in physics. Michael Faraday visualized magnetic forces as "lines of force" stretching through space. Later, James Clerk Maxwell translated Faraday's qualitative ideas into rigorous mathematics, giving birth to classical electromagnetism. 

Simultaneously, mathematicians and physicists studying fluid dynamics needed a way to describe the velocity of water particles in a river. A vector field perfectly models fluid flow: at any point $(x,y)$ in the river, the vector $\mathbf{F}(x,y)$ represents the speed and direction of the water flowing through that exact spot. Thus, vector fields bridge the gap between abstract geometry and the physical realities of gravity, electricity, magnetism, and fluid mechanics.

\subsection{Key Formulas}
\begin{itemize}
    \item \textbf{General 2D Vector Field:} $\mathbf{F}(x, y) = \langle P(x, y), Q(x, y) \rangle$
    \item \textbf{General 3D Vector Field:} $\mathbf{F}(x, y, z) = \langle P(x, y, z), Q(x, y, z), R(x, y, z) \rangle$
    \item \textbf{Magnitude of a Vector Field:} $|\mathbf{F}(x, y)| = \sqrt{P(x,y)^2 + Q(x,y)^2}$
    \item \textbf{Gradient Vector Field:} $\nabla f(x, y) = \left\langle \frac{\partial f}{\partial x}, \frac{\partial f}{\partial y} \right\rangle$
    \item \textbf{Linear Approximation for Position (Euler's Method):} If $\mathbf{V}(x,y)$ is a velocity field, the new position after a small time step $\Delta t$ is approximately:
    \[ \text{Position}_{\text{new}} \approx \text{Position}_{\text{old}} + \mathbf{V}(x, y) \Delta t \]
\end{itemize}

\subsection{Prerequisites}
To master this topic, you must be comfortable with:
\begin{enumerate}
    \item \textbf{Vector Algebra:} Magnitudes, scalar multiplication, dot products, and writing vectors in standard basis form ($\mathbf{i}, \mathbf{j}, \mathbf{k}$) vs. component form ($\langle x, y, z \rangle$).
    \item \textbf{Partial Derivatives:} Applying the chain rule, product rule, and quotient rule with respect to a single variable while treating other variables as constants.
    \item \textbf{Graphing Basics:} Identifying lines, circles, and quadrants in the Cartesian plane to intuitively sketch how arrows should point.
\end{enumerate}

% ==============================================================================
% PART 2: Detailed Homework Solutions
% ==============================================================================
\section{Part 2: Detailed Homework Solutions}

\subsection*{Problem 1}
\textbf{Question:} Sketch the vector field $\mathbf{F}(x, y) = \frac{y\mathbf{i} + x\mathbf{j}}{\sqrt{x^2 + y^2}}$. Identify the correct plot.

\textbf{Solution:}
Let's analyze the properties of this vector field by checking its magnitude and direction.
\begin{enumerate}
    \item \textbf{Magnitude:} Let's find the length of any vector in this field.
    \[ |\mathbf{F}(x,y)| = \sqrt{ \left(\frac{y}{\sqrt{x^2+y^2}}\right)^2 + \left(\frac{x}{\sqrt{x^2+y^2}}\right)^2 } = \sqrt{\frac{y^2}{x^2+y^2} + \frac{x^2}{x^2+y^2}} = \sqrt{\frac{x^2+y^2}{x^2+y^2}} = \sqrt{1} = 1 \]
    Except at the origin $(0,0)$, every vector has a length exactly equal to $1$. The correct plot must show arrows of constant length.
    \item \textbf{Direction:} Test test points on the axes:
    \begin{itemize}
        \item At $(1, 0)$ (positive x-axis): $\mathbf{F}(1, 0) = \frac{0\mathbf{i} + 1\mathbf{j}}{\sqrt{1}} = \langle 0, 1 \rangle$. The arrow points \textit{straight up}.
        \item At $(0, 1)$ (positive y-axis): $\mathbf{F}(0, 1) = \frac{1\mathbf{i} + 0\mathbf{j}}{\sqrt{1}} = \langle 1, 0 \rangle$. The arrow points \textit{straight right}.
        \item At $(1, 1)$ (Quadrant I): $\mathbf{F}(1, 1) = \frac{1\mathbf{i} + 1\mathbf{j}}{\sqrt{2}} = \left\langle \frac{1}{\sqrt{2}}, \frac{1}{\sqrt{2}} \right\rangle$. The arrow points \textit{up and right}.
        \item At $(-1, 1)$ (Quadrant II): $\mathbf{F}(-1, 1) = \frac{1\mathbf{i} - 1\mathbf{j}}{\sqrt{2}} = \left\langle \frac{1}{\sqrt{2}}, -\frac{1}{\sqrt{2}} \right\rangle$. The arrow points \textit{down and right}.
    \end{itemize}
\end{enumerate}
By observing these test points, the correct graph must have arrows forming a "saddle" pattern, flowing inward along $y=-x$ and outward along $y=x$, with all arrows having the same length.

\textbf{Final Answer:} The correct plot matches the saddle-like flow with uniform arrow lengths (typically the first or second graph in such WebAssign sets, matched exactly via the test vectors above).

\subsection*{Problem 2}
\textbf{Question:} Sketch the vector field $\mathbf{F}(x, y, z) = \mathbf{i}$. Evaluate it at 8 specific points.

\textbf{Solution:}
The vector field is given by $\mathbf{F}(x, y, z) = 1\mathbf{i} + 0\mathbf{j} + 0\mathbf{k} = \langle 1, 0, 0 \rangle$.
Notice that this function has no $x$, $y$, or $z$ variables in its output. This means it is a \textbf{constant vector field}. At every single coordinate in 3D space, the vector is identical: it points exactly 1 unit in the positive x-direction.

Evaluating the field at the requested points:
\begin{align*}
    \text{At } (-2, -2, -2) &\implies \mathbf{F} = \langle 1, 0, 0 \rangle \\
    \text{At } (-2, -2, 2)  &\implies \mathbf{F} = \langle 1, 0, 0 \rangle \\
    \text{At } (-2, 2, -2)  &\implies \mathbf{F} = \langle 1, 0, 0 \rangle \\
    \text{At } (-2, 2, 2)   &\implies \mathbf{F} = \langle 1, 0, 0 \rangle \\
    \text{At } (2, -2, -2)  &\implies \mathbf{F} = \langle 1, 0, 0 \rangle \\
    \text{At } (2, -2, 2)   &\implies \mathbf{F} = \langle 1, 0, 0 \rangle \\
    \text{At } (2, 2, -2)   &\implies \mathbf{F} = \langle 1, 0, 0 \rangle \\
    \text{At } (2, 2, 2)    &\implies \mathbf{F} = \langle 1, 0, 0 \rangle
\end{align*}

\textbf{Final Answer:} Every single input box should be filled with the vector $\langle 1, 0, 0 \rangle$.

\subsection*{Problem 3}
\textbf{Question:} Match the vector field $\mathbf{F}(x, y) = \langle x, -y \rangle$ with the correct plot.

\textbf{Solution:}
Test specific points to determine the behavior of the arrows.
\begin{itemize}
    \item Positive x-axis $(x, 0)$: $\mathbf{F}(x, 0) = \langle x, 0 \rangle$. Arrows point outwards along the x-axis, increasing in length as $|x|$ increases.
    \item Positive y-axis $(0, y)$: $\mathbf{F}(0, y) = \langle 0, -y \rangle$. Arrows point \textit{downward} towards the origin along the y-axis.
    \item Quadrant I $(1, 1)$: $\mathbf{F}(1, 1) = \langle 1, -1 \rangle$. Arrows point down and right.
\end{itemize}
This flow pushes particles toward the x-axis and away from the y-axis, creating a classic "saddle" or hyperbolic flow.

\textbf{Final Answer:} Select the plot where arrows point inwards along the vertical axis and outwards along the horizontal axis.

\subsection*{Problem 4}
\textbf{Question:} Match the vector field $\mathbf{F}(x, y) = \langle y, 2x \rangle$ with the correct plot.

\textbf{Solution:}
Test axes to determine the flow:
\begin{itemize}
    \item On the x-axis $(x, 0)$: $\mathbf{F}(x, 0) = \langle 0, 2x \rangle$. These arrows point purely vertically. For $x > 0$, they point straight up. For $x < 0$, they point straight down.
    \item On the y-axis $(0, y)$: $\mathbf{F}(0, y) = \langle y, 0 \rangle$. These arrows point purely horizontally. For $y > 0$, they point right. For $y < 0$, they point left.
    \item Quadrant I $(1, 1)$: $\mathbf{F}(1, 1) = \langle 1, 2 \rangle$. The vertical component is twice as strong as the horizontal component, meaning arrows are steeply pointing up and right.
\end{itemize}

\textbf{Final Answer:} Select the plot that has vertical vectors on the x-axis (pointing up on the right, down on the left) and horizontal vectors on the y-axis.

\subsection*{Problem 5}
\textbf{Question:} Match the vector field on $\mathbb{R}^3$, $\mathbf{F}(x, y, z) = x\mathbf{i} + y\mathbf{j} + 3\mathbf{k}$.

\textbf{Solution:}
Analyze the components separately.
1. The $x$ and $y$ components are $\langle x, y \rangle$. In the xy-plane, this represents a field that points radially outward from the origin, growing larger as you move farther away.
2. The $z$ component is a constant $3$. This means every single arrow in the entire space is tilted "upward" (in the positive z-direction) by a constant amount of 3.

\textbf{Final Answer:} The correct 3D plot will show arrows exploding outward away from the z-axis, while simultaneously all pointing diagonally upward at a constant z-rate.

\subsection*{Problem 6}
\textbf{Question:} Find the partial derivatives and gradient vector field $\nabla f$ of $f(x, y) = 6y \sin(xy)$.

\textbf{Solution:}
To find the gradient field, we must calculate the partial derivative with respect to $x$ ($f_x$) and the partial derivative with respect to $y$ ($f_y$).

\textbf{Step 1: Find $f_x(x, y)$}
Treat $y$ as a constant. Use the chain rule for $\sin(xy)$.
\[ f_x(x, y) = \frac{\partial}{\partial x} [6y \sin(xy)] \]
\[ f_x(x, y) = 6y \cdot \cos(xy) \cdot \frac{\partial}{\partial x}(xy) \]
\[ f_x(x, y) = 6y \cdot \cos(xy) \cdot (y) = 6y^2 \cos(xy) \]

\textbf{Step 2: Find $f_y(x, y)$}
Treat $x$ as a constant. Here, we must use the \textbf{Product Rule} because we have $y$ multiplied by $\sin(xy)$, both containing the variable $y$.
\[ f_y(x, y) = \frac{\partial}{\partial y} [6y \cdot \sin(xy)] \]
\[ f_y(x, y) = \left[ \frac{\partial}{\partial y}(6y) \right] \sin(xy) + 6y \left[ \frac{\partial}{\partial y}(\sin(xy)) \right] \]
\[ f_y(x, y) = 6 \sin(xy) + 6y \cdot \cos(xy) \cdot (x) \]
\[ f_y(x, y) = 6 \sin(xy) + 6xy \cos(xy) \]

\textbf{Step 3: Construct $\nabla f$}
The gradient vector field is simply the vector containing both partial derivatives.
\[ \nabla f(x, y) = \langle f_x(x, y), f_y(x, y) \rangle \]

\textbf{Final Answers:}
$f_x(x, y) = 6y^2 \cos(xy)$ \\
$f_y(x, y) = 6\sin(xy) + 6xy \cos(xy)$ \\
$\nabla f(x, y) = \langle 6y^2 \cos(xy), 6\sin(xy) + 6xy \cos(xy) \rangle$

\subsection*{Problem 7}
\textbf{Question:} Find the gradient vector field $\nabla f$ of $f(s, t) = \sqrt{7s + 4t}$.

\textbf{Solution:}
First, rewrite the function with a fractional exponent for easier differentiation:
\[ f(s, t) = (7s + 4t)^{1/2} \]

\textbf{Step 1: Find $f_s(s, t)$}
Treat $t$ as a constant. Use the generalized power rule (chain rule).
\[ f_s(s, t) = \frac{1}{2}(7s + 4t)^{-1/2} \cdot \frac{\partial}{\partial s}(7s + 4t) \]
\[ f_s(s, t) = \frac{1}{2}(7s + 4t)^{-1/2} \cdot (7) = \frac{7}{2\sqrt{7s + 4t}} \]

\textbf{Step 2: Find $f_t(s, t)$}
Treat $s$ as a constant.
\[ f_t(s, t) = \frac{1}{2}(7s + 4t)^{-1/2} \cdot \frac{\partial}{\partial t}(7s + 4t) \]
\[ f_t(s, t) = \frac{1}{2}(7s + 4t)^{-1/2} \cdot (4) = \frac{4}{2\sqrt{7s + 4t}} = \frac{2}{\sqrt{7s + 4t}} \]

\textbf{Step 3: Construct $\nabla f$}
\[ \nabla f(s, t) = \langle f_s(s, t), f_t(s, t) \rangle \]

\textbf{Final Answer:} 
$\nabla f(s, t) = \left\langle \frac{7}{2\sqrt{7s + 4t}}, \frac{2}{\sqrt{7s + 4t}} \right\rangle$

\subsection*{Problem 8}
\textbf{Question:} Match the function $f(x,y) = (x+y)^2$ with the correct gradient vector field plot.

\textbf{Solution:}
First, compute the gradient vector field $\nabla f(x, y)$.
\[ f_x(x, y) = 2(x+y)^1 \cdot (1) = 2(x+y) \]
\[ f_y(x, y) = 2(x+y)^1 \cdot (1) = 2(x+y) \]
\[ \nabla f(x, y) = \langle 2(x+y), 2(x+y) \rangle \]

Analyze the behavior of this field:
1. Since the x-component and y-component are always perfectly equal to each other, every single vector must point at a 45-degree angle (parallel to the line $y=x$ or $y=-x$).
2. Consider the line $y = -x$. If we plug this into the gradient, we get $\nabla f = \langle 2(0), 2(0) \rangle = \langle 0, 0 \rangle$. The arrows should shrink to zero length exactly along the line $y = -x$.
3. As you move away from the line $y=-x$, the magnitude $2(x+y)$ gets larger, so the arrows should grow in length.

\textbf{Final Answer:} Select the plot where all vectors are perfectly diagonal (slope of 1), and there is an empty diagonal band (zero vector length) passing through the origin along $y = -x$.

\subsection*{Problem 9}
\textbf{Question:} A particle moves in a velocity field $\mathbf{V}(x, y) = \langle x^2, x + y^2 \rangle$. If it is at position $(6, 5)$ at time $t = 6$, estimate its location at time $t = 6.01$.

\textbf{Solution:}
This is a real-world application of a vector field acting as a rate of change. We can use linear approximation (Euler's Method).
The formula relating displacement, velocity, and time is:
\[ \Delta \mathbf{r} \approx \mathbf{V}(x_0, y_0) \Delta t \]
where $\Delta \mathbf{r} = \langle \Delta x, \Delta y \rangle$ is the change in position.

\textbf{Step 1: Find the change in time $\Delta t$.}
\[ \Delta t = t_{\text{final}} - t_{\text{initial}} = 6.01 - 6 = 0.01 \]

\textbf{Step 2: Evaluate the velocity vector at the current position.}
Plug $(x, y) = (6, 5)$ into the vector field $\mathbf{V}(x, y)$:
\[ \mathbf{V}(6, 5) = \langle (6)^2, 6 + (5)^2 \rangle \]
\[ \mathbf{V}(6, 5) = \langle 36, 6 + 25 \rangle = \langle 36, 31 \rangle \]
This means the particle is moving at a speed of $36$ units/sec in the x-direction and $31$ units/sec in the y-direction.

\textbf{Step 3: Calculate the displacement.}
\[ \Delta \mathbf{r} \approx \langle 36, 31 \rangle \times (0.01) = \langle 0.36, 0.31 \rangle \]

\textbf{Step 4: Find the final position.}
\[ \text{New Position} = \text{Old Position} + \text{Displacement} \]
\[ \text{New Position} = (6, 5) + (0.36, 0.31) = (6.36, 5.31) \]

\textbf{Final Answer:} The estimated location is $(6.36, 5.31)$.

\subsection*{Problem 10}
\textbf{Question:} Fill in the blanks analyzing the vector field $\mathbf{F}(x,y) = -y\mathbf{i} + x\mathbf{j}$ on $\mathbb{R}^2$.

\textbf{Solution:}
The field in component form is $\mathbf{F}(x,y) = \langle -y, x \rangle$. Let's trace the blanks from the image.

\textbf{Intro text blanks:}
\begin{itemize}
    \item We know the basis vector $\mathbf{j}$ is defined as $\langle 0, 1 \rangle$.
    \item "Since $\mathbf{F}(1,0) = \mathbf{j}$, we draw the vector \textbf{$\langle 0, 1 \rangle$} starting at the point $(1,0)$."
    \item "Since $\mathbf{F}(0,1) = -\mathbf{i}$, we draw the vector $\langle -1, 0 \rangle$ with starting point \textbf{$(0, 1)$}."
\end{itemize}

\textbf{Table Blanks:}
We evaluate $\mathbf{F}(x,y) = \langle -y, x \rangle$ for given coordinates, or work backwards if the output is given.
\begin{itemize}
    \item Row 2: $\mathbf{F}(2, 2) = \langle -2, 2 \rangle$. (Left box)
    \item Row 4: We are given the input $(0,1)$ yielding $\langle -1,0 \rangle$. The right side has a blank input $(x,y)$ that yields the output vector $\langle 1, 0 \rangle$.
    Set $\langle -y, x \rangle = \langle 1, 0 \rangle$. This means $-y = 1 \implies y = -1$, and $x = 0$. So the input is \textbf{$(0, -1)$}.
    \item Row 6: We have a blank input for $(x,y)$ on the left yielding output $\mathbf{F}(x,y) = \langle -3, 0 \rangle$.
    Set $\langle -y, x \rangle = \langle -3, 0 \rangle \implies -y = -3 \implies y = 3$, and $x = 0$. The input is \textbf{$(0, 3)$}.
\end{itemize}

\textbf{Dot Product Check:}
To confirm arrows are tangent to circles centered at the origin, check the dot product between the position vector $\mathbf{x} = \langle x, y \rangle$ and the field vector $\mathbf{F}(\mathbf{x}) = \langle -y, x \rangle$:
\[ \mathbf{x} \cdot \mathbf{F}(\mathbf{x}) = (x\mathbf{i} + y\mathbf{j}) \cdot (-y\mathbf{i} + x\mathbf{j}) \]
Multiply the components:
\[ = \textbf{-xy} + \textbf{yx} \]
\[ = \textbf{0} \]
Since the dot product is zero, the position vector (pointing radially out) is strictly perpendicular to the flow vector. Therefore, the flow is purely rotational (circular).

% ==============================================================================
% PART 3: Analysis of Problems and Techniques
% ==============================================================================
\section{Part 3: In-Depth Analysis of Problems and Techniques}

\subsection{A) Problem Types and General Approach}
The homework categorizes neatly into four distinct problem types:
\begin{enumerate}
    \item \textbf{Vector Field Sketching / Matching (Probs 1, 3, 4, 5, 8):} The strategy is to heavily rely on "test points." The easiest points are on the x-axis (where $y=0$) and y-axis (where $x=0$). By evaluating the vector field equation at $(1,0)$ and $(0,1)$, you can immediately eliminate 3 out of 4 multiple-choice graphs.
    \item \textbf{Direct Evaluation (Prob 2, 10):} This simply requires plugging coordinates into the component functions $P(x,y,z)$ and $Q(x,y,z)$. 
    \item \textbf{Gradient Vector Fields (Probs 6, 7):} The strategy is entirely calculus-based. You are given a scalar function $f(x,y)$ and must compute the partial derivatives $f_x$ and $f_y$ to form the vector $\langle f_x, f_y \rangle$.
    \item \textbf{Linear Approximation / Kinematics (Prob 9):} This treats the vector field as a physical entity (velocity). The core strategy relies on the relationship $\text{Distance} = \text{Velocity} \times \text{Time}$, adapted for vectors: $\Delta \mathbf{r} = \mathbf{V} \cdot \Delta t$.
\end{enumerate}

\subsection{B) Key Algebraic and Calculus Manipulations}
\begin{itemize}
    \item \textbf{Magnitude Simplification Trick:} In Problem 1, verifying uniform arrow length required algebraically isolating the magnitude: $|\mathbf{F}| = \sqrt{\frac{y^2}{x^2+y^2} + \frac{x^2}{x^2+y^2}}$. Recognizing that the numerator naturally combines into a common denominator $x^2+y^2$ allowed the expression to collapse to $\sqrt{1} = 1$.
    \item \textbf{Product Rule with Partial Derivatives:} In Problem 6, evaluating $\frac{\partial}{\partial y} [6y \sin(xy)]$ was a common trap. Students often forget the product rule because they are hyper-focused on treating $x$ as a constant. Because $y$ exists in \textit{both} terms ($6y$ and $\sin(xy)$), the product rule $(fg)' = f'g + fg'$ is strictly required.
    \item \textbf{Dot Product for Orthogonality:} Problem 10 utilized the geometric property of the dot product: $\mathbf{u} \cdot \mathbf{v} = 0 \iff \mathbf{u} \perp \mathbf{v}$. By proving that $\langle x,y \rangle \cdot \langle -y, x \rangle = -xy + yx = 0$, we mathematically proved the field perfectly models rotational fluid flow.
\end{itemize}

% ==============================================================================
% PART 4: Cheatsheet and Tips
% ==============================================================================
\section{Part 4: "Cheatsheet" and Tips for Success}

\subsection{Summary of Most Important Formulas}
\begin{itemize}
    \item $\nabla f(x, y) = \langle f_x, f_y \rangle$ \quad (Gradient Field)
    \item $|\mathbf{F}(x,y)| = \sqrt{P(x,y)^2 + Q(x,y)^2}$ \quad (Arrow Length)
    \item Rotational Field Prototype: $\mathbf{F}(x,y) = \langle -y, x \rangle$ (Counterclockwise circular flow)
    \item Radial Field Prototype: $\mathbf{F}(x,y) = \langle x, y \rangle$ (Exploding outwards from origin)
\end{itemize}

\subsection{Tips, Tricks, and Shortcuts}
\begin{itemize}
    \item \textbf{The "Axis Trick" for Graph Matching:} \textit{Never} try to compute vectors in the middle of a quadrant first. Always plug in $x=0$ (the y-axis) and $y=0$ (the x-axis). If you plug in $x=0$ and the resulting vector field has no $\mathbf{j}$ component, it means all arrows crossing the y-axis are perfectly horizontal. Look for the graph that matches this!
    \item \textbf{The "Component Zero" Trick:} Look at the field $\mathbf{F} = \langle P, Q \rangle$. Set $P = 0$ to find where the arrows are completely vertical. Set $Q = 0$ to find where the arrows are completely horizontal. 
    \item \textbf{Gradient Check:} If you have a contour map (level curves) of a function, the gradient vector field arrows will \textit{always} point perpendicular to the level curves, pointing towards the highest contour values.
\end{itemize}

\subsection{Common Traps}
\begin{itemize}
    \item \textbf{Mixing up components:} When calculating $\nabla f = \langle f_x, f_y \rangle$, writing $f_y$ first by accident. Always keep strict alphabetical order in your vector components.
    \item \textbf{Vectors vs. Points:} A point $(x,y)$ dictates \textit{where} the tail of the arrow sits. The vector $\langle P, Q \rangle$ dictates the \textit{geometry of the arrow} itself. Never confuse the two sets of coordinates.
\end{itemize}

% ==============================================================================
% PART 5: Conceptual Synthesis
% ==============================================================================
\section{Part 5: Conceptual Synthesis and The "Big Picture"}

\subsection{A) Thematic Connections}
The central theme of this topic is \textbf{"Assigning multiple rates of change to every point in space."} 
In Calc I, the derivative $f'(x)$ assigned a single slope to every x-value. In Calc III, a vector field assigns a directional magnitude to every coordinate. Think of a vector field as the spatial equivalent of a derivative function. Just as $f'(x)$ tells you how a 2D graph behaves at a point, a vector field $\mathbf{F}$ tells you how an object in a physical space will be pushed or pulled at any exact location.

\subsection{B) Forward and Backward Links}
\textbf{Backward Links:} Vector fields directly utilize your knowledge from earlier in Calc III of partial derivatives and basis vectors ($\mathbf{i}, \mathbf{j}, \mathbf{k}$). The gradient vector $\nabla f$, which previously evaluated the direction of steepest ascent at one specific point, is now zoomed out to show the steepest ascent direction everywhere simultaneously.

\textbf{Forward Links:} This is arguably the most crucial foundation for the rest of Calculus III. Almost everything you do from here on out involves vector fields. You will learn \textbf{Line Integrals} ($\int \mathbf{F} \cdot d\mathbf{r}$), which calculate the "work" done by a vector field pushing an object along a path. You will also learn \textbf{Green's Theorem, Stokes' Theorem, and the Divergence Theorem}, which relate the macroscopic flow and rotation of a vector field to specific boundary integrals.

% ==============================================================================
% PART 6: Real-World Application (Finance/Economics)
% ==============================================================================
\section{Part 6: Real-World Application and Modeling}

While vector fields are famously used in physics (electromagnetism and fluid dynamics), they are deeply embedded in modern quantitative finance and macroeconomics.

\subsection{A) Concrete Scenarios in Finance and Economics}
\textbf{1. Quantitative Finance: Markowitz Portfolio Optimization and Gradient Descent.}
In modern portfolio theory, an algorithmic trader wishes to minimize the variance (risk) of a portfolio containing multiple assets while achieving a target return. If the portfolio contains two assets with weights $w_1$ and $w_2$, the risk is a surface function $\sigma^2(w_1, w_2)$. To write an algorithm that automatically adjusts asset weights to minimize risk, quants use the gradient vector field $\nabla \sigma^2$. The algorithm "flows" backwards along the vector field vectors ($-\nabla \sigma^2$) to slide down the risk surface to the absolute minimum.

\textbf{2. Macroeconomics: State-Space Phase Portraits.}
Central banks model the economy as a dynamical system. Imagine a 2D space where the x-axis is the Inflation Rate ($\pi$) and the y-axis is the Interest Rate ($r$). Based on the Taylor Rule and Phillips Curve, economists create a vector field $\mathbf{F}(\pi, r) = \langle \dot{\pi}, \dot{r} \rangle$. By looking at this vector field, economists can instantly see if the economy is spiraling out of control (arrows pushing outward) or stabilizing towards a healthy equilibrium (arrows flowing inward to a target point).

\subsection{B) Model Problem Setup: The Risk Gradient Field}
Let's model the Quantitative Finance scenario.
Suppose a quant constructs a portfolio of two stocks: Stock A and Stock B.
\begin{itemize}
    \item Let $w_1$ be the fraction of money invested in Stock A.
    \item Let $w_2$ be the fraction of money invested in Stock B.
\end{itemize}
Based on historical data, the covariance and variance of the stocks create a Risk Function $R(w_1, w_2)$ modeled by:
\[ R(w_1, w_2) = 4w_1^2 + 9w_2^2 + 2w_1w_2 \]
The algorithmic trading system needs to know which direction to shift the weights to increase the risk the fastest (so the algorithm can do the exact opposite). This requires finding the gradient vector field of the Risk Function.

\textbf{The Math Setup:}
The system must compute the vector field:
\[ \mathbf{F}(w_1, w_2) = \nabla R(w_1, w_2) = \left\langle \frac{\partial R}{\partial w_1}, \frac{\partial R}{\partial w_2} \right\rangle \]
Taking the partial derivatives yields:
\[ \mathbf{F}(w_1, w_2) = \langle 8w_1 + 2w_2, 18w_2 + 2w_1 \rangle \]
If the current portfolio weights are $w_1 = 0.6$ and $w_2 = 0.4$, plugging those coordinates into the vector field $\mathbf{F}$ outputs an exact directional arrow. The algorithm will then step in the direction of $-\mathbf{F}(0.6, 0.4)$ to effectively lower the risk.

% ==============================================================================
% PART 7: Common Variations and Untested Concepts
% ==============================================================================
\section{Part 7: Common Variations and Untested Concepts}

A major concept introduced in Section 16.1 that was conspicuously absent from your homework set is the test for \textbf{Conservative Vector Fields}.

\subsection{What is a Conservative Vector Field?}
A vector field $\mathbf{F}$ is called "conservative" if it is the gradient of some scalar function $f$ (meaning $\mathbf{F} = \nabla f$). Such a function $f$ is called the \textbf{potential function} for $\mathbf{F}$. In physics, gravity is a conservative vector field, meaning the work done moving an object depends only on the start and end point, not the path taken.

\subsection{The Cross-Partial Test}
If $\mathbf{F}(x, y) = \langle P(x, y), Q(x, y) \rangle$ is defined on a simply-connected open region, it is conservative \textit{if and only if}:
\[ \frac{\partial P}{\partial y} = \frac{\partial Q}{\partial x} \]
If these partial derivatives are equal, the field is conservative.

\textbf{Worked Example:}
Determine whether $\mathbf{F}(x, y) = \langle x^2y, xy^2 \rangle$ is a conservative vector field.
\begin{enumerate}
    \item Identify $P$ and $Q$: $P(x,y) = x^2y$ and $Q(x,y) = xy^2$.
    \item Differentiate $P$ with respect to $y$: $\frac{\partial P}{\partial y} = x^2(1) = x^2$.
    \item Differentiate $Q$ with respect to $x$: $\frac{\partial Q}{\partial x} = (1)y^2 = y^2$.
    \item Compare: Since $x^2 \neq y^2$, the cross-partials are \textbf{not} equal. 
\end{enumerate}
\textit{Conclusion:} This vector field is \textbf{not} conservative. It is impossible to find a scalar function $f(x,y)$ such that $\nabla f = \mathbf{F}$.

% ==============================================================================
% PART 8: Advanced Diagnostic Testing: "Find the Flaw"
% ==============================================================================
\section{Part 8: Advanced Diagnostic Testing: "Find the Flaw"}
Below are 5 problems with subtle flaws. Find the error, explain it, and fix it!

\subsection*{Problem 1: 3D Gradient Vector Field}
\textbf{Question:} Find the gradient vector field $\nabla f$ for $f(x,y,z) = x^2 y z^3$. \\
\textbf{Flawed Solution:}
1. $f_x = 2x y z^3$ \\
2. $f_y = x^2 z^3$ \\
3. $\nabla f(x,y,z) = \langle 2x y z^3, x^2 z^3 \rangle$ \\
\textbf{Your Task:} Find the flaw.

\vspace{0.2cm}
\noindent \textit{\textbf{Correction:}} The flaw is that the student entirely ignored the z-component $f_z$; a 3D scalar function must produce a 3D gradient vector field. \\
\textit{Correct Step:} Calculate $f_z = x^2 y (3z^2) = 3x^2 y z^2$. \\
\textit{Correct Final Answer:} $\nabla f(x,y,z) = \langle 2x y z^3, x^2 z^3, 3x^2 y z^2 \rangle$.

\subsection*{Problem 2: Proving a Gradient Field}
\textbf{Question:} Determine if $\mathbf{F}(x,y) = \langle y, x \rangle$ is the gradient field of $f(x,y) = xy$. \\
\textbf{Flawed Solution:}
1. Find $f_x$: $\frac{\partial}{\partial x}(xy) = y$. \\
2. Find $f_y$: $\frac{\partial}{\partial y}(xy) = x$. \\
3. Because $f_x$ is $y$ and $f_y$ is $x$, the gradient vector is $\langle f_y, f_x \rangle = \langle x, y \rangle$. \\
4. Since $\langle x, y \rangle \neq \langle y, x \rangle$, it is NOT the gradient field. \\
\textbf{Your Task:} Find the flaw.

\vspace{0.2cm}
\noindent \textit{\textbf{Correction:}} The flaw is in step 3 where the student constructed the gradient vector backwards as $\langle f_y, f_x \rangle$ instead of the standard $\langle f_x, f_y \rangle$. \\
\textit{Correct Step:} The definition of the gradient is $\nabla f = \langle f_x, f_y \rangle = \langle y, x \rangle$. \\
\textit{Correct Final Answer:} Yes, $\mathbf{F}(x,y)$ \textit{is} the gradient field of $f(x,y) = xy$.

\subsection*{Problem 3: Linear Approximation (Euler's Method)}
\textbf{Question:} A particle moves in velocity field $\mathbf{V}(x,y) = \langle y, -x \rangle$. It is at $(1, 0)$ at $t=0$. Estimate its position at $t=0.1$. \\
\textbf{Flawed Solution:}
1. Evaluate velocity at start: $\mathbf{V}(1, 0) = \langle 0, -1 \rangle$. \\
2. Multiply by time step $\Delta t = 0.1$: $\text{Displacement} = \langle 0, -0.1 \rangle$. \\
3. New position = $(0, 0) + \langle 0, -0.1 \rangle = (0, -0.1)$. \\
\textbf{Your Task:} Find the flaw.

\vspace{0.2cm}
\noindent \textit{\textbf{Correction:}} The flaw is in step 3; the student added the displacement to the origin $(0,0)$ instead of the particle's actual starting position $(1,0)$. \\
\textit{Correct Step:} $\text{New position} = \text{Old position} + \text{Displacement} = (1, 0) + (0, -0.1)$. \\
\textit{Correct Final Answer:} $(1, -0.1)$.

\subsection*{Problem 4: Dot Product Geometry}
\textbf{Question:} Show that the vector field $\mathbf{F}(x,y) = \langle x, y \rangle$ points parallel to the position vector $\mathbf{r} = \langle x, y \rangle$ by computing the dot product. \\
\textbf{Flawed Solution:}
1. $\mathbf{r} \cdot \mathbf{F} = \langle x, y \rangle \cdot \langle x, y \rangle$. \\
2. $\mathbf{r} \cdot \mathbf{F} = \langle x^2, y^2 \rangle$. \\
3. Because the components are positive squares, the vectors are parallel. \\
\textbf{Your Task:} Find the flaw.

\vspace{0.2cm}
\noindent \textit{\textbf{Correction:}} The flaw is a fundamental misunderstanding of the dot product; the dot product of two vectors yields a scalar (a number), not a vector. \\
\textit{Correct Step:} $\langle x, y \rangle \cdot \langle x, y \rangle = x^2 + y^2$. \\
\textit{Correct Final Answer:} The dot product is the scalar $x^2 + y^2$. (To prove they are parallel, one should technically show $\mathbf{F} = 1\mathbf{r}$, not use the dot product, but the dot product math itself was fatally flawed).

\subsection*{Problem 5: Product Rule in Partial Derivatives}
\textbf{Question:} Find $f_x$ for $f(x,y) = x^2 \cos(xy)$. \\
\textbf{Flawed Solution:}
1. Treat $y$ as a constant. \\
2. Derivative of $x^2$ is $2x$. \\
3. Derivative of $\cos(xy)$ with respect to $x$ is $-y \sin(xy)$. \\
4. Multiply them together: $f_x = -2xy \sin(xy)$. \\
\textbf{Your Task:} Find the flaw.

\vspace{0.2cm}
\noindent \textit{\textbf{Correction:}} The flaw is taking the derivatives of two terms containing $x$ separately and multiplying them, completely ignoring the mandatory Product Rule. \\
\textit{Correct Step:} $f_x = \frac{\partial}{\partial x}(x^2) \cdot \cos(xy) + x^2 \cdot \frac{\partial}{\partial x}(\cos(xy))$. \\
\textit{Correct Final Answer:} $f_x = 2x \cos(xy) - x^2 y \sin(xy)$.

\end{document}